\documentclass[12pt]{article}
\usepackage[utf8]{inputenc}
\usepackage[T1]{fontenc}
\usepackage{hyperref}
\title{Hukumina Language Resources for K-12 Curriculum Development}
\author{Compiled by Claude Sonnet 5.5 (claude-sonnet-5-5) for Andr\'as Kornai}
\date{}
\begin{document}
\maketitle

\section{Introduction}
Hukumina (ISO 639-3 huw) is an extinct Austronesian language of northwest Buru, Maluku, Indonesia. Sources differ on when it died out: one speaker is reported in 1989, and ASJP lists it as extinct since 1995. Search results also say Indonesia's education ministry listed it among 11 extinct languages in 2024. Only scholarly and encyclopedic material was found.

\section{Primary Educational Resources}
None found. No teaching materials or revitalization programme appeared in the search results.

\section{Linguistic and Reference Materials}
\begin{itemize}
\item Overview: \url{https://en.wikipedia.org/wiki/Hukumina_language}
\item C.\ Grimes, ``Digging for the Roots of Language Death in Eastern Indonesia: The Cases of Kayeli and Hukumina'' (ANU repository; only the abstract and record page were seen): \url{https://openresearch-repository.anu.edu.au/handle/1885/39691?mode=full}
\item \emph{Endangered Languages of Austronesia} (Oxford University Press), chapter listing: \url{https://academic.oup.com/book/27194/chapter/196640810}
\item ASJP wordlist database entry: \url{https://asjp.clld.org/languages/HUKUMINA}
\end{itemize}

\section{Religious and Community Resources}
None found.

\section{Digital and Audio Resources}
No audio or digital corpus was found beyond the ASJP wordlist entry above.

\section{Curriculum Development Context}
With no known speakers and little documentation, a heritage curriculum could only cover the history of the Hukumina people and of language loss on Buru, drawing on the scholarly sources above. Any language content would be limited to the short wordlist data. The Buru community should be consulted before anything is built.
\end{document}